\subsubsection*{Question a) Amplitude d'oscillation du bloc}

\begin{enumerate}[label=\arabic*)]
    \item \textbf{Variables connues et recherchées :}
    \begin{itemize}[label= -]
        \item Énergie mécanique : $E = \qty{\mhsqSixEm}{J}$
        \item Masse du bloc : $m = \qty{\mhsqSixMasse}{g} = \qty{\fpeval{\mhsqSixMasse / 1000}}{kg}$
        \item Constante de rappel : $k = \qty{\mhsqSixK}{N/m}$
        \item Variable recherchée : Amplitude $A \text{ (m)}$
    \end{itemize}

    \item \textbf{Équations pertinentes :}
    \[
    E = \frac{1}{2} k A^2
    \]

    \item \textbf{Isolation de la variable recherchée :}
    \[
    A^2 = \frac{2E}{k} \quad \implies \quad A = \sqrt{\frac{2E}{k}}
    \]

    \item \textbf{Résolution numérique :}
    \[
    A = \sqrt{\frac{2 \cdot \mhsqSixEm}{\mhsqSixK}} = \qty{\fpeval{round(sqrt(2 * \mhsqSixEm / \mhsqSixK), 4)}}{m}
    \]

    \item \textbf{Réponse :}
    L'amplitude d'oscillation du bloc est de \textbf{\qty{\fpeval{round(sqrt(2 * \mhsqSixEm / \mhsqSixK), 3)}}{m}} (soit \textbf{\qty{\fpeval{round(sqrt(2 * \mhsqSixEm / \mhsqSixK) * 100, 2)}}{cm}}).
\end{enumerate}


% ------------------------------------------------------------------------------
% QUESTION B
% ------------------------------------------------------------------------------
\subsubsection*{Question b) Module de la vitesse maximale}

\begin{enumerate}[label=\arabic*)]
    \item \textbf{Variables connues et recherchées :}
    \begin{itemize}[label= -]
        \item Énergie mécanique : $E = \qty{\mhsqSixEm}{J}$
        \item Masse du bloc : $m = \qty{\fpeval{\mhsqSixMasse / 1000}}{kg}$
        \item Variable recherchée : Vitesse maximale $v_{\max} \text{ (m/s)}$
    \end{itemize}

    \item \textbf{Équations pertinentes :}
    \[
    E = \frac{1}{2} m v_{\max}^2
    \]

    \item \textbf{Isolation de la variable recherchée :}
    \[
    v_{\max}^2 = \frac{2E}{m} \quad \implies \quad v_{\max} = \sqrt{\frac{2E}{m}}
    \]

    \item \textbf{Résolution numérique :}
    \[
    v_{\max} = \sqrt{\frac{2 \cdot \mhsqSixEm}{\fpeval{\mhsqSixMasse / 1000}}} = \qty{\fpeval{round(sqrt(2 * \mhsqSixEm / (\mhsqSixMasse / 1000)), 2)}}{m/s}
    \]

    \item \textbf{Réponse :}
    Le module de la vitesse maximale atteinte par le bloc est de \textbf{\qty{\fpeval{round(sqrt(2 * \mhsqSixEm / (\mhsqSixMasse / 1000)), 2)}}{m/s}}.
\end{enumerate}


% ------------------------------------------------------------------------------
% QUESTION C
% ------------------------------------------------------------------------------
\subsubsection*{Question c) Position lorsque la vitesse est de \qty{\mhsqSixVit}{m/s}}

\begin{enumerate}[label=\arabic*)]
    \item \textbf{Variables connues et recherchées :}
    \begin{itemize}[label= -]
        \item Énergie mécanique : $E = \qty{\mhsqSixEm}{J}$
        \item Masse du bloc : $m = \qty{\fpeval{\mhsqSixMasse / 1000}}{kg}$
        \item Constante de rappel : $k = \qty{\mhsqSixK}{N/m}$
        \item Vitesse instantanée : $v = \qty{\mhsqSixVit}{m/s}$
        \item Variable recherchée : Position $x \text{ (m)}$
    \end{itemize}

    \item \textbf{Équations pertinentes :}
    \[
    E = K + U = \frac{1}{2} m v^2 + \frac{1}{2} k x^2
    \]

    \item \textbf{Isolation de la variable recherchée :}
    \[
    \frac{1}{2} k x^2 = E - \frac{1}{2} m v^2 \quad \implies \quad x^2 = \frac{2E - m v^2}{k} \quad \implies \quad x = \pm \sqrt{\frac{2E - m v^2}{k}}
    \]

    \item \textbf{Résolution numérique :}
    \[
    K = \frac{1}{2} \cdot \left(\fpeval{\mhsqSixMasse / 1000}\right) \cdot (\mhsqSixVit)^2 = \num[scientific-notation=true, round-mode=figures, round-precision=3]{\fpeval{0.5 * (\mhsqSixMasse / 1000) * (\mhsqSixVit)^2}} \text{ J}
    \]
    \[
    x = \pm \sqrt{\frac{2 \cdot \mhsqSixEm - (\fpeval{\mhsqSixMasse / 1000}) \cdot (\mhsqSixVit)^2}{\mhsqSixK}} = \pm \qty{\fpeval{round(sqrt((2 * \mhsqSixEm - (\mhsqSixMasse / 1000) * (\mhsqSixVit)^2) / \mhsqSixK), 4)}}{m}
    \]

    \item \textbf{Réponse :}
    Lorsque la vitesse du bloc est de \qty{\mhsqSixVit}{m/s}, sa position se trouve à \textbf{$\pm$ \qty{\fpeval{round(sqrt((2 * \mhsqSixEm - (\mhsqSixMasse / 1000) * (\mhsqSixVit)^2) / \mhsqSixK), 3)}}{m}} (soit \textbf{$\pm$ \qty{\fpeval{round(sqrt((2 * \mhsqSixEm - (\mhsqSixMasse / 1000) * (\mhsqSixVit)^2) / \mhsqSixK) * 100, 2)}}{cm}}) par rapport à sa position d'équilibre.
\end{enumerate}